\documentclass[11pt]{article}
\usepackage[margin=1in]{geometry}
\usepackage{amsmath,amssymb,amsfonts}
\usepackage{array}
\usepackage{graphicx}
\usepackage{microtype}
\usepackage[hidelinks]{hyperref}

\newcommand{\R}{\mathbb{R}}
\newcommand{\Obs}{\mathcal{O}}
\newcommand{\Act}{\mathcal{A}}
\newcommand{\St}{\mathcal{S}}
\newcommand{\Cset}{\mathcal{C}}
\newcommand{\Dset}{\mathcal{D}}
\newcommand{\qref}{q^{\mathrm{ref}}}
\renewcommand{\vec}[1]{\mathbf{#1}}
\newcommand{\clip}{\operatorname{clip}}
\newcommand{\diag}{\operatorname{diag}}

\title{\textbf{Recovering Muscle Weakness from Kinematics Alone\\ in a Musculoskeletal Arm Simulation:\\ State/Action Spaces, Reference Generation, and Identification}}
\author{myoArm identifiability study --- technical report}
\date{2026-06-11}

\begin{document}
\maketitle

\begin{abstract}
We study \emph{sim-internal identifiability}: in a 63-muscle MuJoCo arm (myoArm), we scale per-muscle maximum isometric force ($s_F$) and optimal fiber length ($s_L$), let a reinforcement-learning policy \emph{softly} track human elevation references, and then recover $(s_F,s_L)$ from the resulting motion. We give formal definitions of the (i) Markov decision process state, observation, and action spaces; (ii) the reward; (iii) the policy/value output spaces and PPO objective; (iv) two generative models for the reference trajectories (functional PCA / ProMP and a VAE) with their input/output spaces; and (v) the time-series identification regressor. Using \textbf{kinematics only} (no EMG), a temporal convolutional regressor recovers shoulder/elbow $s_F$ at $R^2\!=\!0.81$; unlocking the wrist with the full movement adds wrist-muscle identifiability ($R^2\!\approx\!0.8$) at the cost of shoulder/elbow accuracy.
\end{abstract}

\section{Problem formulation}
Let the plant be the myoArm model with generalized coordinates $\vec q\in\R^{n_q}$ ($n_q{=}38$: $27$ independent $+\,11$ scapular-rhythm coordinates coupled by equality constraints), generalized velocities $\dot{\vec q}\in\R^{n_v}$, and $n_\mu{=}63$ Hill-type muscle actuators with activation states $\vec a^{\mathrm{act}}\in[0,1]^{n_\mu}$. A \emph{perturbation} (muscle weakness) is a pair of channel-wise scale vectors
\begin{equation}
\boldsymbol\theta=(\vec s_F,\vec s_L)\in\Theta:=[0.05,1]^{|\Cset|}\times[0.70,1.20]^{|\Cset|},
\end{equation}
applied to $|\Cset|$ \emph{perturbation channels} (groups of muscles). For channel $c$ acting on muscle $i$, the MuJoCo gain parameters are scaled as
\begin{equation}
\text{gain}_{i,2}\!\leftarrow\! s_F^{c}\,\text{gain}^{0}_{i,2}\ \ (\text{gain}_2{=}F_{\max}),\qquad
\text{gain}_{i,0:2}\!\leftarrow\!\text{gain}^{0}_{i,0:2}/s_L^{c}\ \ (\text{operating length range}).
\end{equation}
The goal is to learn an estimator $\widehat{\boldsymbol\theta}=g(\text{observed motion})$ and report per-channel recovery quality $R^2_c$. The motion is produced by a policy that softly tracks a reference trajectory $\qref:[0,H]\to\R^{|\Dset|}$ over the tracked degrees of freedom $\Dset$.

\section{Markov decision process}
The control MDP is $\mathcal M=(\St,\Obs,\Act,P,r,\gamma)$ at control rate $50$\,Hz ($\Delta t{=}0.02$\,s, physics substeps $F{=}10$ at $500$\,Hz), horizon $H$ ($175$ for elevation-rise, $200$ for the full sequence).

\subsection{State space $\St$}
The (Markov) state augments the physical state with the episode-fixed context:
\begin{equation}
\St=\underbrace{\R^{n_q}}_{\vec q}\times\underbrace{\R^{n_v}}_{\dot{\vec q}}\times\underbrace{[0,1]^{n_\mu}}_{\vec a^{\mathrm{act}}}\times\underbrace{\Theta}_{\text{weakness}}\times\underbrace{\Lambda}_{\text{latent}}\times\underbrace{\{T_1,T_2\}}_{\text{task}}\times\underbrace{\{0,\dots,H\}}_{\text{step }t},
\end{equation}
where $\Lambda\subset\R^{5}$ collects the movement-variation latent $\boldsymbol\ell=(T,\mathrm{peak},\mathrm{skew},\mathrm{plane},\mathrm{elbow})$ (Sec.~\ref{sec:ref}). $\boldsymbol\theta,\boldsymbol\ell,\text{task}$ are drawn at reset and held fixed within an episode; $\qref$ is a deterministic function of $(\text{task},\boldsymbol\ell)$.

\subsection{Observation space $\Obs$}
The policy does not see $\St$ directly but an observation $\vec o_t=\phi(s_t)\in\Obs\subset\R^{d_o}$, $d_o=76$ (wrist-unlocked mode: $96$). With tracked DoF $\Dset=\{\text{elv},\text{plane},\text{elbow},\text{prosup}\}$ ($|\Dset|{=}4$) and free DoF $\text{rot}$, and writing $\vec a\in\Act$ for the previous action, $\phi$ concatenates
\begin{align}
\vec o_t=\Big[\ &\tfrac{1}{\pi}\,\vec q_{\Dset\cup\{\text{rot}\}}\ (5),\quad \tfrac{1}{5}\,\dot{\vec q}_{\Dset\cup\{\text{rot}\}}\ (5),\quad \vec a^{\mathrm{act}}_{\Act}\ (26),\nonumber\\
&\tfrac{1}{\pi}\qref_{\Dset}(t)\ (4),\quad \tfrac{1}{\pi}\big(\qref_{\Dset}(t){-}\vec q_{\Dset}\big)\ (4),\quad \tfrac{1}{\pi}\qref_{\Dset}(t{+}5)\ (4),\nonumber\\
&\big[\varphi_t,\ \mathbb 1_{\text{hold}}\big]\ (2),\quad \tilde{\boldsymbol\ell}\ (5),\quad [\vec s_F,\vec s_L]\ (2|\Cset|{=}20),\quad \mathbb 1[\text{task}{=}T_2]\ (1)\ \Big],\label{eq:obs}
\end{align}
with phase $\varphi_t=t/H$, $\mathbb 1_{\text{hold}}=\mathbb 1[t\ge T/\Delta t]$, and normalized latent $\tilde{\boldsymbol\ell}=(T/3,\mathrm{peak}/120,\mathrm{skew}/2,\mathrm{plane}/45,\mathrm{elbow}/130)$. Dimension check: $5{+}5{+}26{+}4{+}4{+}4{+}2{+}5{+}20{+}1=76$. A running-statistics affine normalizer $\nu$ (VecNormalize, clipped to $[-10,10]$) is applied before the network: $\tilde{\vec o}_t=\clip\!\big((\vec o_t-\boldsymbol\mu_o)/\boldsymbol\sigma_o,\,{-}10,10\big)$. The block $[\vec s_F,\vec s_L]$ encodes the (privileged) weakness, justified by the \emph{adapted chronic} scenario: the CNS has re-adapted to fixed muscle parameters, so motor commands reflect $\boldsymbol\theta$.

\subsection{Action space $\Act$ and dynamics}
\begin{equation}
\Act=[-1,1]^{|\Act|},\qquad |\Act|=26\ \text{(shoulder 15 + elbow 9 + pronators PT,PQ)},
\end{equation}
mapped to muscle excitation $u\in[0,1]^{|\Act|}$ with signal-dependent motor noise (Harris--Wolpert):
\begin{equation}
u_j=\tfrac12(a_j+1),\qquad u_j\leftarrow\clip\!\big(u_j+\eta_j,0,1\big),\quad \eta_j\sim\mathcal N\!\big(0,(\sigma_{\mathrm{sd}}u_j+\sigma_0)^2\big),
\end{equation}
$\sigma_{\mathrm{sd}}{=}0.1$, $\sigma_0{=}0.01$. Non-action muscles take $u{=}0$. MuJoCo maps excitation to activation through first-order dynamics ($\tau_{\mathrm{act}}{=}10$\,ms, $\tau_{\mathrm{deact}}{=}40$\,ms) and force through a Hill force--length--velocity model; transition $P$ integrates $F$ physics substeps.

\subsection{Reward}
\begin{equation}
r(s_t,\vec a_t)=\mathrm{TASK}_t\cdot\mathrm{QUALITY}_t+w_e\,\mathrm{EFFORT}_t+\mathrm{pen}^{\mathrm{rot}}_t,
\end{equation}
\begin{align}
\mathrm{TASK}_t&=\sum_{i\in\Dset} w_i\,\exp\!\Big(-k_i^{\mathrm{out}}\big[\max(0,\,|q_i-\qref_i|-b_i)\big]^2\Big),\\
\mathrm{QUALITY}_t&=\exp\!\Big(-\big(c_J\,\overline{\ddot q^2}+c_D\,\overline{(\vec a_t-\vec a_{t-1})^2}+c_S\,\mathbb 1_{\text{hold}}\,\overline{\dot q^2}\big)\Big),\\
\mathrm{EFFORT}_t&=-\tfrac{1}{|\Act|}\textstyle\sum_{j}(a^{\mathrm{act}}_j)^{p},\qquad
\mathrm{pen}^{\mathrm{rot}}_t=-w_r\big[\max(0,|q_{\text{rot}}|-50^\circ)\big]^2,
\end{align}
where $\overline{(\cdot)}$ is the mean over $\Dset$. The two-sided \emph{soft band}: deviations within $b_i$ incur \emph{no} penalty (allowing self-paced variation), beyond $b_i$ a Gaussian falloff; this preserves weakness as deviation rather than absorbing it. Coefficients (locked): $w=(0.15,0.15,0.25,0.20)$, $b=(10,15,12,20)^\circ$, $k^{\mathrm{out}}=(30,18,22,12)$, $w_e{=}0.05$, $p{=}2$, $c_J{=}5\!\times\!10^{-6}$, $c_D{=}0.5$, $c_S{=}1.0$, $w_r{=}0.5$. Episode terminates (with $r\!-\!1$) if $|q_{\text{elv}}{-}\qref_{\text{elv}}|>50^\circ$, $|q_{\text{elbow}}{-}\qref_{\text{elbow}}|>60^\circ$, or $\max|\dot q|>40$.

\section{Policy, value, and PPO}
\subsection{Output spaces}
The stochastic policy maps observations to a distribution over actions,
\begin{equation}
\pi_\psi:\Obs\to\mathcal P(\Act),\qquad \pi_\psi(\cdot\mid\vec o)=\mathcal N\!\big(\boldsymbol\mu_\psi(\vec o),\,\diag(e^{2\boldsymbol\sigma_\psi})\big),
\end{equation}
so the \emph{policy output space} is $\R^{|\Act|}\times\R_{>0}^{|\Act|}$ (mean $\boldsymbol\mu_\psi(\vec o)\in\R^{26}$ and a state-independent log-std $\boldsymbol\sigma_\psi\in\R^{26}$). The value head outputs $V_\psi:\Obs\to\R$. Both $\boldsymbol\mu_\psi$ and $V_\psi$ are MLPs with hidden layers $[256,256]$ and $\tanh$ activations (separate networks).

\subsection{Objective}
With generalized advantage $\hat A_t=\sum_{l\ge0}(\gamma\lambda)^l\delta_{t+l}$, $\delta_t=r_t+\gamma V_\psi(\vec o_{t+1})-V_\psi(\vec o_t)$, and ratio $\rho_t(\psi)=\pi_\psi(\vec a_t\mid\vec o_t)/\pi_{\psi_{\mathrm{old}}}(\vec a_t\mid\vec o_t)$, PPO minimizes
\begin{equation}
\mathcal L(\psi)=-\,\mathbb E_t\!\Big[\min\!\big(\rho_t\hat A_t,\ \clip(\rho_t,1{-}\epsilon,1{+}\epsilon)\hat A_t\big)\Big]
+c_v\,\mathbb E_t\!\big[(V_\psi(\vec o_t)-\hat R_t)^2\big]-c_e\,\mathbb E_t\!\big[\mathcal H(\pi_\psi(\cdot\mid\vec o_t))\big].
\end{equation}
Hyperparameters: $\epsilon{=}0.2$, learning rate $3\!\times\!10^{-4}$, rollout $n_{\text{steps}}\!\times\!n_{\text{env}}=512\!\times\!20$, minibatch $4096$, $10$ epochs/update, $\gamma{=}0.99$, $\lambda{=}0.95$, $c_v{=}0.5$, $c_e{=}0$, gradient-norm clip $0.5$. Curriculum/warm-start: M1 healthy fixed (T1) $\to$ M2 latent distribution $+$ T2 $\to$ M3 perturbation $k{=}1\!\to\!k\!\le\!3$; weights are warm-started across stages ($\Obs$ dimension held fixed).

\section{Reference (imitation target) generation}\label{sec:ref}
The reference $\qref$ is generated from human data (KIMHu, 20 subjects, Kinect-V2 3D joints). Per-repetition multichannel rise (or full-cycle) curves are derived \emph{from 3D positions only} and resampled to $N$ phase points, yielding a dataset
\begin{equation}
\mathcal X=\{\vec X^{(m)}\in\R^{N\times d}\}_{m=1}^{M},\quad (N,d)=(50,4)\ \text{[rise]}\ \text{or}\ (80,6)\ \text{[full+wrist]},\ M\approx200,
\end{equation}
with channels (elevation, elbow, plane, pronation $[+$ wrist flexion, deviation$]$). After per-channel standardization, $\vec x^{(m)}=\operatorname{vec}\big((\vec X^{(m)}-\boldsymbol\mu)\oslash\boldsymbol\sigma\big)\in\R^{Nd}$.

\subsection{Functional PCA / ProMP (adopted)}
\textbf{Input space} $\R^{Nd}$; \textbf{output} a sampled trajectory in $\R^{N\times d}$ (then time-warped to $\R^{H\times|\Dset|}$). Center $\vec x_c^{(m)}=\vec x^{(m)}-\bar{\vec x}$, take the SVD $\mathbf X_c=\mathbf U\mathbf S\mathbf V^\top$, eigenvalues $\lambda_k=S_{kk}^2/(M{-}1)$, principal functions $\boldsymbol\phi_k=\mathbf V_{:,k}$, and retain
\begin{equation}
K=\min\Big\{K:\ \textstyle\sum_{k\le K}\lambda_k\big/\sum_k\lambda_k\ge0.95\Big\}\quad(K=13\text{--}24\ \text{empirically}).
\end{equation}
Standardized scores $\hat w^{(m)}_k=\langle\vec x_c^{(m)},\boldsymbol\phi_k\rangle/\sqrt{\lambda_k}$. \textbf{Generation} (kernel density in score space; bootstraps a real score vector with Gaussian jitter, preserving non-Gaussian structure and the empirical covariance):
\begin{equation}
\hat{\vec z}=\hat{\vec w}^{(j)}+\boldsymbol\varepsilon,\ \ j\!\sim\!\mathrm{Unif}\{1{:}M\},\ \boldsymbol\varepsilon\!\sim\!\mathcal N(0,h^2 I),\ h{=}0.35;\quad
\hat{\vec x}=\bar{\vec x}+\sum_{k=1}^{K}\hat z_k\sqrt{\lambda_k}\,\boldsymbol\phi_k.
\end{equation}
With $\hat{\vec z}\!\sim\!\mathcal N(0,I)$ the generated covariance equals the (truncated) empirical covariance $\sum_k\lambda_k\boldsymbol\phi_k\boldsymbol\phi_k^\top$, so dispersion is preserved (no mode shrinkage). A duration $T\!\sim\!\widehat{\mathrm{Emp}}\{T^{(m)}\}$ time-warps $\hat{\vec X}$ via $\varphi(t)=\min(t/T,1)$ to give $\qref$. Empirically peak $94{\pm}9^\circ$ (real) vs.\ $95{\pm}9^\circ$ (fPCA); marginal Kolmogorov--Smirnov distance $0.07$--$0.14$.

\subsection{Variational autoencoder (baseline)}
\textbf{Encoder} $q_\vartheta:\R^{Nd}\to\mathcal P(\R^{6})$ and \textbf{decoder} $p_\vartheta:\R^{6}\to\R^{Nd}$ (latent space $\R^{6}$):
\begin{equation}
\boldsymbol\mu_z,\log\boldsymbol\sigma_z^2=\mathrm{MLP}_{200\to128\to128\to(6,6)}(\vec x),\quad
\hat{\vec x}=\mathrm{MLP}_{6\to128\to128\to200}(\vec z),\quad \vec z=\boldsymbol\mu_z+\boldsymbol\sigma_z\odot\boldsymbol\varepsilon,
\end{equation}
trained by the $\beta$-ELBO ($\beta{=}0.5$)
\begin{equation}
\mathcal L_{\mathrm{VAE}}=\|\vec x-\hat{\vec x}\|_2^2+\beta\cdot\tfrac12\textstyle\sum_k\big(\sigma_{z,k}^2+\mu_{z,k}^2-1-\log\sigma_{z,k}^2\big),
\end{equation}
Adam ($10^{-3}$, $1500$ epochs). Generation: $\vec z\!\sim\!\mathcal N(0,I)\!\to\!p_\vartheta$. The KL term plus small $M$ induce \emph{under-dispersion} (peak s.d.\ $9^\circ\!\to\!4^\circ$) and partial posterior collapse ($3/6$ active dims); fPCA is therefore preferred.

\section{Identification regressor}
\subsection{Input/output spaces}
From $N\!\approx\!16$--$20$k policy rollouts (motor noise on; deterministic policy mean), each episode yields a \emph{kinematics-only} multivariate time series and a label:
\begin{equation}
g:\ \underbrace{\R^{T\times C}}_{\text{kinematics }\vec Z}\times\underbrace{\R^{6}}_{\text{known covariates }\boldsymbol\ell}\ \longrightarrow\ \underbrace{\R^{2|\Cset|}}_{\widehat{\boldsymbol\theta}=[\hat{\vec s}_F;\hat{\vec s}_L]},
\end{equation}
with $C=2(|\Dset|{+}1)+|\Dset|$ channels (positions and velocities of tracked DoF $+$ rotation, plus tracking errors): $C{=}14$ (base) or $20$ (wrist). \emph{No muscle activation / EMG is used.} Variable lengths are zero-padded to $T_{\max}$ with mask $\vec m\in\{0,1\}^{T}$.

\subsection{Architecture (dilated TCN with masked pooling)}
\begin{align}
\vec H^{(0)}&=\mathrm{ReLU}\,\mathrm{GN}\,\mathrm{Conv1d}_{k=7,s=2}(\vec Z^\top),\qquad
\vec H^{(l)}=\mathrm{ReLU}\,\mathrm{GN}\,\mathrm{Conv1d}_{k=5,\,\mathrm{dil}=2^{l}}(\vec H^{(l-1)}),\ l{=}1,2,3,\\
\vec g&=\Big[\ \tfrac{\sum_t m_t\vec H_t}{\sum_t m_t}\ ;\ \max_{t:m_t=1}\vec H_t\ ;\ \boldsymbol\ell\ \Big],\qquad
\widehat{\boldsymbol\theta}=\mathrm{MLP}_{\cdot\to128\to128\to 2|\Cset|}(\vec g),
\end{align}
channel widths $64\!\to\!64\!\to\!128\!\to\!128$, GroupNorm, masked average$+$maximum temporal pooling. Loss with severity reweighting $w_n=1+2\max_c(1-s_F^{c,n})$ and Huber:
\begin{equation}
\mathcal L=\frac{1}{|\mathcal B|}\sum_{n\in\mathcal B}w_n\cdot\frac{1}{2|\Cset|}\sum_o \mathrm{Huber}\big(\hat\theta^{(n)}_o-\theta^{(n)}_o\big),
\end{equation}
Adam ($10^{-3}$, weight decay $10^{-5}$) with cosine schedule, batch $256$, $100$--$120$ epochs, $85/15$ hold-out (best-val). Evaluation: per-channel $R^2$ and MAE.

\section{Results}
Kinematics-only recovery of $s_F$ (Fmax scale), fPCA references with motor noise (Table~\ref{tab:res}). The temporal model is essential: a snapshot (mean/max summary) MLP attains only $R^2\!\approx\!0.36$ on the same kinematics, versus $0.81$ for the TCN --- the weakness signal lives in the \emph{temporal structure} (rise slope, deviation onset).

\begin{table}[h]\centering
\begin{tabular}{lccc}
\hline\hline
Setup & input ch.\ $C$ & shoulder/elbow $R^2(s_F)$ & wrist $R^2(s_F)$\\
\hline
rise, T1$+$T2, wrist locked & 14 & \textbf{0.81} & --\\
full-sequence $+$ wrist, T1 & 20 & 0.60 & 0.81\\
full-sequence $+$ wrist, T1$+$T2 & 20 & 0.55 & 0.74\\
\hline\hline
\end{tabular}
\caption{Per-configuration kinematics-only identifiability. Unlocking the wrist with the full movement adds wrist-muscle identification but dilutes shoulder/elbow accuracy, because the elevation-rise phase (muscles maximally loaded against gravity) is the most informative segment.}
\label{tab:res}
\end{table}

\noindent\textbf{Per-channel} (rise, best setup): A\_lowcuff $.94$, B\_latadd $.88$, SUPSP/BIClong/TRIlong $.86$--$.87$, DELT1 $.84$, DELT2 $.79$, DELT3 $.75$, CORB/PECM1 $.62$--$.65$. The length scale $s_L$ is harder from kinematics alone ($R^2\!\approx\!0.56$).

\paragraph{Conclusions.} (i) Muscle \emph{force} weakness of shoulder/elbow muscles is strongly recoverable from movement kinematics alone ($R^2{=}0.81$), without EMG. (ii) Reference generation by fPCA/ProMP reproduces the empirical trajectory distribution faithfully (no under-dispersion), outperforming a VAE. (iii) The soft-band reward is the mechanism that turns weakness into an identifiable signal (mild $\!\to\!$ activation compensation, severe $\!\to\!$ kinematic deviation). (iv) There is a protocol trade-off: one movement cannot simultaneously maximize shoulder/elbow and wrist identifiability.

\end{document}
